\chapter{Conclusion}
\label{chapter:conclusion}

This thesis investigated two complementary questions: whether replacing the sequential
ConvGRU fusion in SOFI-MISRGRU with a Transformer's parallel attention fusion yields
better super-resolution quality or lower latency. Secondly, whether MISRGRU can be 
restructured into a streaming architecture that produces one super-resolved output per 
incoming frame, enabling continuous real-time SR video rather than periodic batch outputs.
Both questions were evaluated on a synthetic fluorescence microscopy dataset spanning
twelve SNR conditions (three blinking speeds $\times$ four photon levels), using the
resolution-scaled Pearson correlation (RSP) and decorrelation analysis as primary metrics
to match the evaluation protocol of the foundational RESURF work~\cite{resurf2024}.

\section*{TR-MISR versus MISRGRU (RQ1)}

The TR-MISR architecture, adapted from the satellite SR domain to second-order SOFI upscaling
($2H{-}7$ output size), produces SR reconstructions that are competitive with the 20-frame
MISRGRU foundational model.
On fast-blinking and high-SNR conditions the two architectures perform similarly in RSP;
on slow-blinking and low-SNR conditions, TR-MISR holds its quality while MISRGRU degrades
more sharply and exhibits hallucinated filaments.
The advantage of the attention mechanism is therefore most visible precisely where it
matters most for live-cell imaging: difficult acquisitions with sparse blinking events and
low photon counts.

TR-MISR does not, however, outperform MISRGRU on raw inference latency.
At $N{=}20$ frames the lightest TR-MISR configuration (1 attention head, 1 Transformer
layer; 1h1d) takes 262\,ms per reconstruction compared to 170\,ms for MISRGRU on the same
hardware.
An ablation across four Transformer sizes (1h1d through 8h6d) showed that RSP saturates
immediately: quality does not improve meaningfully beyond 1h1d despite a $3\times$ increase
in parameters.
Latency, by contrast, grows steeply. From 88\,ms to 209\,ms at $N{=}8$ going from 1h1d
to 4h3d.
The encoder and decoder are the architectural bottleneck for quality, not the size 
of the fusion module.
Future TR-MISR improvements should target deeper residual encoders or deformable
convolutional decoders rather than larger Transformers.

\section*{Frame generalisation of TR-MISR}

The most practically useful property of the Transformer fusion module is its
length-invariance at inference time.
Because the cross-attention softmax redistributes weight over all input tokens rather than
accumulating magnitude, the fused representation fed to the decoder remains at a stable
scale regardless of how many frames are provided.
% Measured decoder-input L2-norms confirm this: TR-MISR's fusion output varies by less than
% 3\% as frame count increases from 8 to 30, whereas MISRGRU's summative GAP layer produces
% norms that grow from 671 to 5\,786 over the same range for the 20-frame foundational model.

This stability means that a TR-MISR model trained on 8 frames can be evaluated on 20 to 40
frames at test time without retraining or architectural changes, if the memory allows it.
Providing 20 frames to an 8-frame-trained model matches or exceeds the RSP of the
20-frame MISRGRU baseline.
A counterintuitive but robust finding is that an 8-frame-trained TR-MISR outperforms a
20-frame-trained TR-MISR when both are evaluated on 20-frame inputs.
Training under the tighter 8-frame budget forces the attention module to assign
discriminative weights to every available frame; training on 20 frames risks over-saturating
the cross-attention on redundant repeated blinking observations of the same fluorophores.
The practical consequence is that TR-MISR can be trained on short sequences, reducing
compute cost and training time, while the researcher retains the option of feeding a
longer sequence at inference for improved quality.

\section*{Streaming MISRGRU and real-time deployment (RQ2)}

The streaming MISRGRU variant, which emits one super-resolved frame for every input frame
by carrying the ConvGRU hidden state across the entire acquisition without resetting,
addresses the throughput bottleneck of the batch model.
Batch MISRGRU produces one SR output per $N$ acquired frames; the streaming model produces
one SR output per frame.

At $N{=}20$ frames on a 512$\times$512 field of view, the streaming model takes
11.29\,ms per SR output on an NVIDIA RTX A5000 which is a $15\times$ improvement over the batch
model's 170.29\,ms per output at the same frame count.
At 20 accumulated frames the streaming model's RSP converges to the batch model's value,
and it continues to improve beyond that point as the GRU accumulates more context across
the acquisition a capability the batch model lacks entirely. S

Training the streaming model required truncated backpropagation through time (TBPTT) to
keep memory tractable for long sequences.
Supervising only the last $n$ frames of each TBPTT chunk was critical: supervising every
frame from the first step prevented the GRU from learning to defer quality improvements to
later frames.
An ablation across TBPTT window length, loss function, GRU width, and decoder type
determined the final configuration: LLN35 supervision window, Fourier-only loss, a single
48-channel ConvGRU, and a transposed-convolution decoder.
% The Fourier-only loss was preferred because L1 supervision over the background-dominated
% pixel distribution suppresses fine structural detail; the frequency-domain loss penalises
% missing high-frequency structure directly.

\section*{Motion reconstruction (RQ3)}

The static streaming model was further extended to handle inter-frame sample motion, which
is unavoidable in live-cell imaging.
Two strategies were evaluated: direct finetuning on motion-augmented data, and the addition
of a Deformable Phase-space Alignment (DPA) module that corrects the GRU hidden state
into the coordinate frame of the current input before the recurrent update.
The final DPA model is effectively motion-invariant: at fast/high-SNR its reconstruction
fidelity holds whether the sample is static or moving up to 5.0\,$\mu$m/s, whereas the
static-trained TR-MISR and MISRGRU baselines collapse at those same speeds.
DPA also outperforms the plain finetune-only variant, with visual inspection confirming less
ghosting and blur from misaligned hidden states.

On static data the only degradation is in slow-low-SNR conditions, where long-range temporal
averaging is most valuable and a model trained to absorb inter-frame displacement can no
longer fully exploit it; here the TRMISR keep a small edge.
For deployment, the unmodified streaming model is preferred when the sample is stationary
or drift is sub-pixel; the DPA variant should be selected whenever the specimen exhibits
visible drift relative to the field of view.

\section*{Summary and outlook}

Taken together, the results confirm that deep-learning SR models can produce
second-order-SOFI-quality reconstructions from as few as 8 fluorescence frames in a
fraction of the time required by classical SOFI.
TR-MISR's attention mechanism provides robust quality under low-SNR conditions and a
practically valuable frame-count generalization property that the GRU cannot match.
The streaming MISRGRU architecture removes the batch acquisition bottleneck entirely,
delivering continuous SR video at 11\,ms per frame or 17\,ms for DPA model, comfortably within the 25\,ms budget
of a 40\,fps acquisition, while improving the reconstruction quality with more frames input.

All quantitative results are obtained on synthetic simulated microtubule data.
Transfer to real experimental microscope data remains an open step before the actual deployment.
% A promising route for domain adaptation is to fine-tune the network on a short sequence
% of blank background frames that capture the specific camera noise and PSF of the target
% microscope.
Additionally, higher-order SOFI targets (third or fourth cumulants) could be incorporated
to push spatial resolution further; the decoder architecture is straightforwardly
extensible to this setting.

The combination of a streaming output model achieving real-time SR video
and a Transformer that robustly aggregates frames under challenging imaging conditions
brings deep-learning super-resolution microscopy a meaningful step closer to routine use
in real-time live-cell biology.
